\cotefile{156-8}
\scanpage{106}
\mspage{97}
% archive sheet number: 106

\uline{Ex 3}
\[
\unsure{\Phi|\varphi}=\{S_{ab},P_b,S_{bc},P_d\},\quad
\text{\unsure{comp. connexes}}\quad \{S_{ab},P_b,S_{bc}\}\text{ et }P_d
\]
\[
\begin{aligned}
\delta\Phi&=\{a,b,c,d\}\\
\partial\Phi&=\{a,c\}\deleted{\unsure{d}}\\
\Phi^{\unsure{\bullet|\circ}}&=\{a,\unsure{c'|c},d\}
\end{aligned}
\qquad\text{\unsure{tous les trois différents}.}
\]

\uline{NB} Dans tous les cas on a.
\[
\deleted{\unsure{\Phi}}
\partial\Phi\underbrace{\subset}_{\substack{\text{\unsure{canoniq.}}\\
\text{(\unsure{n'}}\exists\ \text{plus gd \unsure{ni|ou} \inserted{plus petit} élément \unsure{dans} }L\text{ !)}}}
\Phi^\bullet\subset\delta\Phi
\]

Et \deleted{$\Phi^\bullet$ \illegible{2}} On a
\[
\partial\Phi=\left\{a\in\delta\Phi\ \left|\
\begin{gathered}
\inserted{\exists\,X\in\Phi\text{ avec }P_a<X,\text{ et de}}\\
\deleted{\text{ou bien }P_a\notin\Phi,\text{ mais existe }X\in\Phi\text{ où}}\\
\deleted{P_a<X}\\
\text{\inserted{plus,}\quad\inserted{quand } }P_a\in\Phi,\\
\deleted{\text{ou bien }P_a\in\Phi,\text{ alors }\unsure{\exists!}}\\
\text{cet }X\text{ est \uline{unique}}
\end{gathered}\right.\right\}
\]
\margin{\uline{NB} Dire \uline{que} si $P_a\in\Phi$, il existe au plus un $X\in\Phi$ tel que $P_a<X$.}
\[
\begin{aligned}
={}&\text{ens. des sommets non isolés de }\Phi,\\
&\deleted{\text{\unsure{qui sont dans} }\Phi}\ \inserted{\text{qui, de plus, quand ils appartiennent à}}\\
&\deleted{\text{\unsure{qui ne}}}\ \Phi,\text{ ne sont pas ``intérieurs'' à }\Phi
\end{aligned}
\]
\[
\Phi^\bullet\underset{\text{\unsure{canoniq.}}}{=}\partial\Phi\amalg\Phi_{\mathrm{isol}}
\qquad
\begin{gathered}
\text{où }\Phi_{\mathrm{isol}}\text{ est l'ens}\\
\text{des \unsure{ponctuels} isolés de }\Phi
\end{gathered}
\]

Dans $\unsure{\mathscr{M}|\mathcal{M}}$, on va introduire une relation d'ordre, en écrivant
\[
X\underset{\mathrm{pos}}{<}Y
\left\{
\begin{array}{lll}
\text{a)}&X=P_x,\ Y=P_y&x<y\\
\text{b)}&X=P_x,\ Y=S_{ab}&x\leq a<b\\
\text{c)}&X=S_{ab},\ Y=P_y,&a<b\leq y\\
\text{d)}&X=S_{ab},\ Y=S_{a',b'}&a<b\leq a'<b'
\end{array}\right.
\]
% margin note written slanted
\margin{\uline{NB} l'indice pos dans $\underset{\mathrm{pos}}{<}$ est inutile, si on n'a pas \unsure{bloqué} les signes $<,\leq$ pour la relation d'incidence \unsure{ici} dans $\mathscr{M}$\unsure{..}.}
\[
X\underset{\mathrm{pos}}{<}Y\ \Longleftrightarrow\ :\ X\mathrel{|\circ|}Y\text{ et }
\underbrace{\operatorname{or}(\delta X)}_{\substack{\deleted{\text{\unsure{él.}}}\inserted{\text{plus petit}}\\\text{\unsure{élément}}\\\text{de }\delta X}}
<\underbrace{\operatorname{ex}(\delta Y)}_{\substack{\text{plus gd}\\\text{él. de}\\\delta Y}}
\]
